% CaféScan — IEEE Conference Paper
% Compile with: pdflatex -> bibtex -> pdflatex -> pdflatex
\documentclass[conference]{IEEEtran}

\usepackage[T1]{fontenc}
\usepackage[utf8]{inputenc}
\usepackage{amsmath,amssymb}
\usepackage{graphicx}
\usepackage{booktabs}
\usepackage{array}
\usepackage{multirow}
\usepackage{url}
\usepackage{hyperref}
\usepackage{cite}
\usepackage{xcolor}
\usepackage{microtype}
\usepackage{balance}

\hypersetup{
  colorlinks = true,
  linkcolor  = black,
  citecolor  = black,
  urlcolor   = black
}

% -----------------------------------------------------------------------
\begin{document}

\title{A Comparative Study of Deep Learning Architectures\\
for Automated Coffee Leaf Disease Classification\\
Using Multi-Source Image Datasets}

\author{
  \IEEEauthorblockN{Sebastian Martinez Abarca}
  \IEEEauthorblockA{LEAD University\\
    BCD-6210 Advanced Data Mining -- IC 2026\\
    \textit{Email: sebastian.martinez@ulead.ac.cr}}
  \and
  \IEEEauthorblockN{Dr. Juan Murillo-Morera}
  \IEEEauthorblockA{Universidad LEAD\\
    Course Professor\\
    \textit{Email: jmurillo@ulead.ac.cr}}
}

\maketitle

% -----------------------------------------------------------------------
\begin{abstract}
This paper presents a comprehensive comparative study of four deep learning
architectures—EfficientNet-B0, ResNet-50, Vision Transformer (ViT-Small),
and MobileNetV3-Large—for automated classification of coffee leaf diseases
across five agronomic categories: healthy, coffee rust (\textit{Hemileia
vastatrix}), cercospora leaf spot (\textit{Cercospora coffeicola}), leaf
miner (\textit{Leucoptera coffeella}), and Phoma (\textit{Phoma
costaricensis}). A unified multi-source dataset of 59{,}950 images is
assembled from three public collections (BRACOL, JMuBEN, JMuBEN2) and
partitioned using a stratified split by (class $\times$ dataset source) to
prevent domain leakage. All models are fine-tuned from ImageNet
pre-trained weights with differential learning rates, label-smoothing
cross-entropy, class-weighted loss, and mixed-precision training on GPU.
Hyperparameter optimization is performed with Optuna (TPE sampler and
Hyperband pruner). The Vision Transformer achieves the best performance
with a Macro-F1 of 99.65\% and accuracy of 99.71\% on the independent
test set, followed by EfficientNet-B0 (95.40\% Macro-F1), MobileNetV3
(91.38\%), and ResNet-50 (80.39\%). Grad-CAM visualizations confirm that
CNN models attend to disease-relevant image regions. No overfitting is
detected across any model (val--test gap $<$0.2\%). An interactive
Streamlit web application enables real-time inference and model comparison.
\end{abstract}

\begin{IEEEkeywords}
Coffee Leaf Disease, Deep Learning, Vision Transformer, EfficientNet,
Transfer Learning, Grad-CAM, Optuna, Stratified Split, Plant Pathology,
Image Classification, MobileNetV3, ResNet
\end{IEEEkeywords}

% -----------------------------------------------------------------------
\section{Introduction}

Coffee is one of the most economically significant agricultural commodities
globally, with over 175 million 60-kg bags produced annually and more than
125 million livelihoods depending on its cultivation~\cite{ico2023}. In
Central and South America—primary production regions—foliar diseases
represent a major threat to yield stability. Coffee rust alone caused
losses exceeding USD 3.2 billion across Central America during the
2012--2013 outbreak~\cite{bebber2016}, while cercospora, phoma, and leaf
miner collectively reduce productivity by 20--40\% in endemic areas.

Traditional disease diagnosis relies on expert scouting—a labor-intensive
process demanding agronomic knowledge unavailable in many smallholder
contexts. Recent advances in computer vision and deep learning (DL) offer
a promising path toward automated, field-deployable diagnostic tools. Yet
comparative assessments of modern architectures on multi-source coffee
datasets remain limited, particularly with rigorous dataset construction
protocols that prevent leakage between dataset sources.

This study addresses four research questions:
\begin{itemize}
  \item \textbf{RQ1:} How do convolutional neural network architectures
    compare with Vision Transformers for coffee leaf disease classification?
  \item \textbf{RQ2:} How does multi-source dataset composition affect
    model generalization to an independent test set?
  \item \textbf{RQ3:} Which image regions are most discriminative for each
    disease class, as revealed by Grad-CAM?
  \item \textbf{RQ4:} What is the practical trade-off between accuracy and
    computational efficiency for field deployment?
\end{itemize}

% -----------------------------------------------------------------------
\section{Background}

\subsection{Coffee Foliar Diseases}

Five foliar conditions are studied in this work. \textit{Healthy} leaves
exhibit uniform dark-green coloration without visible lesions. \textit{Coffee
rust} (\textit{Hemileia vastatrix}) manifests as yellow-orange powdery
pustules on the abaxial surface, causing severe defoliation. \textit{Cercospora
leaf spot} (\textit{Cercospora coffeicola}) produces circular lesions with a
grey center and yellow halo, affecting leaves, berries, and stems.
\textit{Leaf miner} (\textit{Leucoptera coffeella}) leaves translucent
serpentine galleries created by larval feeding, reducing photosynthetic
area. \textit{Phoma} (\textit{Phoma costaricensis}) causes dark necrotic
lesions with a chlorotic halo, prevalent in high-altitude humid regions.

\subsection{Transfer Learning and Fine-Tuning}

Transfer learning from ImageNet pre-trained weights~\cite{imagenet}
significantly reduces the labeled data requirements for medical and
agricultural image classification~\cite{mohanty2016}. Differential learning
rates—applying a lower rate to the frozen backbone than to the new
classification head—preserve generalized low-level features while adapting
high-level representations to the target domain~\cite{howard2018}.

\subsection{Evaluation Metrics}

For a multi-class problem with $C$ classes, the macro-averaged F1 score is:
\begin{equation}
  \text{Macro-F1} = \frac{1}{C}\sum_{c=1}^{C}
    \frac{2 \cdot P_c \cdot R_c}{P_c + R_c}
  \label{eq:macrof1}
\end{equation}
where $P_c$ and $R_c$ are precision and recall for class $c$.
Macro-F1 is the primary metric because it weights all classes equally,
critical for an imbalanced five-class dataset.

\subsection{Grad-CAM Interpretability}

Gradient-weighted Class Activation Mapping (Grad-CAM)~\cite{selvaraju2017}
generates a saliency map $L^c_{\text{Grad-CAM}}$ for class $c$ from the
last convolutional feature maps $A^k$ and their gradients:
\begin{equation}
  \alpha^c_k = \frac{1}{Z}\sum_{i}\sum_{j}
    \frac{\partial y^c}{\partial A^k_{ij}}, \quad
  L^c = \text{ReLU}\!\left(\sum_k \alpha^c_k A^k\right)
  \label{eq:gradcam}
\end{equation}
Grad-CAM is applicable to CNN architectures; Vision Transformers require
attention rollout or alternative methods.

% -----------------------------------------------------------------------
\section{Related Work}

Mohanty et al.~\cite{mohanty2016} demonstrated the feasibility of deep
learning for plant disease recognition, achieving 99.35\% accuracy on the
PlantVillage dataset using GoogLeNet. However, PlantVillage images were
captured under controlled laboratory conditions, limiting generalization to
field photography.

Dos Santos Ferreira et al.~\cite{ferreira2017} applied CNN-based
classification to Brazilian coffee leaf images (BRACOL), reporting
accuracies exceeding 95\% on single-dataset splits but without multi-source
evaluation. Their work establishes the BRACOL dataset subsequently used in
our study.

Parraga-Alava et al.~\cite{parraga2019} released the JMuBEN dataset,
focusing on rust, cercospora, and phoma identification under Costa Rican
field conditions. They achieved 90--93\% accuracy using AlexNet and VGG-16
with data augmentation.

More recently, Barbedo~\cite{barbedo2019} surveyed deep learning
applications for plant disease detection, highlighting overfitting risks in
single-dataset studies and advocating for multi-source evaluation
protocols—directly motivating our stratified (class $\times$ source) split.

Our work distinguishes itself by: (1) unifying three datasets into a
59{,}950-image benchmark, (2) evaluating four architectures spanning three
design generations (CNN, efficient CNN, Vision Transformer), (3) applying
Bayesian hyperparameter optimization via Optuna, and (4) providing Grad-CAM
interpretability with explicit layer selection per architecture.

% -----------------------------------------------------------------------
\section{Experimental Design}

\subsection{Dataset Construction}

Three publicly available datasets are unified, as summarized in
Table~\ref{tab:datasets}.

\begin{table}[htbp]
  \centering
  \caption{Source Datasets}
  \label{tab:datasets}
  \begin{tabular}{lllp{2.2cm}}
    \toprule
    \textbf{Dataset} & \textbf{Labels} & \textbf{Classes} & \textbf{Origin} \\
    \midrule
    BRACOL   & CSV file        & rust, cercospora, phoma, healthy & Univ. Fed. de Viçosa, Brazil \\
    JMuBEN   & Folder struct.  & rust, cercospora, phoma          & Univ. of Costa Rica \\
    JMuBEN2  & Folder struct.  & healthy, miner                   & Univ. of Costa Rica \\
    \bottomrule
  \end{tabular}
\end{table}

A critical observation motivates the stratified split: \textit{healthy} and
\textit{miner} samples originate exclusively from JMuBEN2, while
\textit{rust}, \textit{cercospora}, and \textit{phoma} come from BRACOL and
JMuBEN. A na\"ive random split would therefore create systematic
class-source confounding between partitions. We address this by stratifying
on the Cartesian product (class $\times$ dataset source), ensuring each
partition receives proportional representation from every (class, source)
stratum. After filtering one corrupted image, the dataset contains
59{,}950 valid samples distributed as shown in Table~\ref{tab:classdist}.

\begin{table}[htbp]
  \centering
  \caption{Class Distribution After Filtering}
  \label{tab:classdist}
  \begin{tabular}{lrrrr}
    \toprule
    \textbf{Class} & \textbf{Total} & \textbf{Train} & \textbf{Val} & \textbf{Test} \\
    \midrule
    healthy    & 20{,}125 & 14{,}087 & 3{,}019 & 3{,}019 \\
    miner      & 17{,}959 & 12{,}571 & 2{,}694 & 2{,}694 \\
    rust       &  9{,}261 &  6{,}482 & 1{,}389 & 1{,}390 \\
    cercospora &  8{,}222 &  5{,}755 & 1{,}233 & 1{,}234 \\
    phoma      &  4{,}383 &  3{,}068 &   657   &   658   \\
    \midrule
    \textbf{Total} & \textbf{59{,}950} & \textbf{41{,}963} & \textbf{8{,}992} & \textbf{8{,}995} \\
    \bottomrule
  \end{tabular}
\end{table}

\subsection{Image Preprocessing and Augmentation}

All images are resized to $256 \times 256$ and randomly cropped to
$224 \times 224$ at training time. The augmentation pipeline applies random
horizontal and vertical flips, color jitter (brightness, contrast,
saturation $\pm 0.3$; hue $\pm 0.1$), and RandAugment~\cite{cubuk2020}
with two operations of magnitude 9. Validation and test images are
deterministically center-cropped to $224 \times 224$ without stochastic
augmentation. All splits are normalized with ImageNet statistics
($\mu = [0.485, 0.456, 0.406]$, $\sigma = [0.229, 0.224, 0.225]$).

\subsection{Model Architectures}

All models are instantiated via the \texttt{timm} library~\cite{rw2019timm}
from ImageNet pre-trained weights. A custom classification head replaces the
original output layer. Parameter counts and head configurations are detailed
in Table~\ref{tab:models}.

\begin{table}[htbp]
  \centering
  \caption{Model Configurations}
  \label{tab:models}
  \begin{tabular}{lrll}
    \toprule
    \textbf{Model} & \textbf{Params} & \textbf{Features} & \textbf{Head} \\
    \midrule
    EfficientNet-B0 & 5.3M  & 1{,}280 & Dropout(0.3) + Linear \\
    ResNet-50       & 25.6M & 2{,}048 & Dropout(0.4) + Linear \\
    ViT-Small/16    & 22.1M & 384     & LayerNorm + Linear     \\
    MobileNetV3-L   & 5.4M  & 960     & Dropout(0.2) + Linear  \\
    \bottomrule
  \end{tabular}
\end{table}

\textbf{EfficientNet-B0}~\cite{tan2019} applies compound scaling of depth,
width, and resolution via a fixed coefficient. Its MBConv blocks with Squeeze-
and-Excitation attention provide parameter efficiency.

\textbf{ResNet-50}~\cite{he2016} uses bottleneck residual blocks with skip
connections, establishing a strong convolutional baseline with 25.6M
parameters.

\textbf{ViT-Small/16}~\cite{dosovitskiy2021} partitions the $224\times224$
image into $196$ non-overlapping $16\times16$ patches, linearly embeds each
patch, and processes the sequence with 8 multi-head self-attention layers.
The \texttt{[CLS]} token representation feeds the classification head.

\textbf{MobileNetV3-Large}~\cite{howard2019} targets mobile deployment
through inverted residual blocks with Hard-Swish activations and Neural
Architecture Search-optimized layer connectivity.

\subsection{Training Protocol}

The loss function combines label smoothing ($\varepsilon = 0.1$) with
class-frequency weighting. For class $c$ with frequency $f_c$ over $C$
classes, the inverse-frequency weight is:
\begin{equation}
  w_c = \frac{\bar{f}}{f_c}, \quad \bar{f} = \frac{1}{C}\sum_{c=1}^C f_c
  \label{eq:classweight}
\end{equation}

The optimizer is AdamW with weight decay $\lambda = 10^{-4}$. Differential
learning rates assign rate $\eta_b = \eta \cdot \gamma$ to backbone
parameters and $\eta$ to the classification head, where $\gamma = 0.1$.
The learning rate schedule follows a linear warmup over 2 epochs, then
cosine decay to zero:
\begin{equation}
  \eta(t) =
  \begin{cases}
    \dfrac{t}{\tau_w}\,\eta & t < \tau_w \\[6pt]
    \dfrac{\eta}{2}\!\left(1 + \cos\dfrac{\pi(t-\tau_w)}{T-\tau_w}\right) & t \geq \tau_w
  \end{cases}
  \label{eq:schedule}
\end{equation}
where $\tau_w$ is the warmup duration and $T$ is the total number of
gradient steps. Mixed-precision training (FP16 on GPU with GradScaler) is
applied throughout. Gradients are clipped to norm 1.0. Training stops when
the validation Macro-F1 does not improve by $\delta = 10^{-4}$ for 8
consecutive epochs (EarlyStopping), and the best checkpoint is restored.

\subsection{Hyperparameter Optimization}

Optuna~\cite{akiba2019} with the TPE sampler and HyperbandPruner explores
the search space shown in Table~\ref{tab:hpo}. The study targets
maximization of validation Macro-F1. Intermediate epoch results are reported
to enable pruning of unpromising trials, and all studies are persisted in a
SQLite database to allow resumable optimization.

\begin{table}[htbp]
  \centering
  \caption{HPO Search Space}
  \label{tab:hpo}
  \begin{tabular}{lll}
    \toprule
    \textbf{Hyperparameter} & \textbf{Range / Choices} & \textbf{Scale} \\
    \midrule
    Learning rate $\eta$        & $[10^{-5},\;10^{-2}]$   & Log     \\
    Weight decay $\lambda$      & $[10^{-5},\;10^{-2}]$   & Log     \\
    Batch size                  & $\{32, 64, 128\}$        & Categ.  \\
    Label smoothing $\varepsilon$ & $[0.0,\;0.2]$          & Linear  \\
    Backbone LR mult. $\gamma$  & $[0.01,\;0.5]$           & Linear  \\
    Warmup epochs               & $\{1, 2, 3, 4, 5\}$      & Int     \\
    \bottomrule
  \end{tabular}
\end{table}

% -----------------------------------------------------------------------
\section{Results}

\subsection{Overview}

Table~\ref{tab:testresults} reports test-set results for all four models.
The ViT-Small achieves dominant performance across all metrics. All models
pass the overfitting check: the absolute gap between validation and test
Macro-F1 is below 0.20 percentage points for every architecture, confirming
that the stratified split successfully prevented data leakage.

\begin{table}[htbp]
  \centering
  \caption{Test Set Results (8{,}993 images)}
  \label{tab:testresults}
  \begin{tabular}{lcccc}
    \toprule
    \textbf{Model} & \textbf{Acc.} & \textbf{Mac-F1} & \textbf{Wgt-F1} & \textbf{Gap} \\
    \midrule
    ViT-Small       & \textbf{99.71} & \textbf{99.65} & \textbf{99.71} & +0.16 \\
    EfficientNet-B0 & 96.45          & 95.40          & 96.44          & +0.18 \\
    MobileNetV3-L   & 93.34          & 91.38          & 93.32          & +0.06 \\
    ResNet-50       & 85.98          & 80.39          & 85.03          & -0.28 \\
    \bottomrule
  \end{tabular}
  \vspace{2pt}
  \footnotesize{All percentages (\%). Gap = val Macro-F1 $-$ test Macro-F1.}
\end{table}

\subsection{Per-Class Analysis}

Table~\ref{tab:perclass} presents per-class F1 scores on the test set,
revealing that \textit{phoma} is the hardest class for CNN architectures.
ResNet-50 achieves only 0.45 F1 for phoma, likely due to phenotypic
similarity between advanced phoma lesions and necrotic cercospora stages.

\begin{table}[htbp]
  \centering
  \caption{Per-Class F1 Score — Test Set}
  \label{tab:perclass}
  \begin{tabular}{lcccc}
    \toprule
    \textbf{Class} & \textbf{ViT} & \textbf{EffNet} & \textbf{MNV3} & \textbf{ResNet} \\
    \midrule
    healthy    & 1.00 & 1.00 & 1.00 & 0.99 \\
    rust       & 1.00 & 0.92 & 0.83 & 0.94 \\
    cercospora & 0.99 & 0.94 & 0.86 & 0.81 \\
    miner      & 1.00 & 0.97 & 0.95 & 0.82 \\
    phoma      & 1.00 & 0.94 & 0.94 & 0.45 \\
    \bottomrule
  \end{tabular}
\end{table}

\subsection{ViT-Small Detailed Results}

The ViT-Small confusion matrix reveals near-perfect classification across
all classes. The only notable error is 8 images of \textit{cercospora}
misclassified as \textit{phoma} (0.7\% of the 1{,}174 cercospora test
samples), which is attributable to the phenotypic overlap between advanced
cercospora lesions and phoma necrosis.

Across all four CNN models, \textit{rust} shows systematically reduced
recall (0.75--0.89), likely because early-stage rust pustules can resemble
pollen deposits or water spots in low-resolution images. ViT correctly
classifies all rust samples, suggesting its global attention mechanism
captures the characteristic spatial distribution of pustules across the
abaxial surface—a feature that local convolutional receptive fields may
not integrate as effectively.

\subsection{EfficientNet-B0 Training Dynamics}

EfficientNet-B0 required 50 epochs to converge (versus 27--38 for the
other three models), with a steady monotonic improvement of Macro-F1 from
0.24 at epoch 1 to 0.96 at epoch 50 without triggering EarlyStopping.
This slow convergence reflects the compound-scaled architecture's need for
more gradient updates to adapt its depth--width--resolution balance to the
new domain. The final test performance (95.40\% Macro-F1) exceeds
MobileNetV3 despite comparable parameter counts, demonstrating that
compound scaling is more effective than depth-only scaling for this task.

\subsection{Overfitting Analysis}

Table~\ref{tab:overfitting} compares validation and test Macro-F1 for all
models. All gaps fall below 0.30 percentage points, confirming the absence
of overfitting. ResNet-50 is the only model where the test F1 marginally
exceeds the validation F1 ($\Delta = -0.28$), indicating that the
validation EarlyStopping signal was slightly conservative.

\begin{table}[htbp]
  \centering
  \caption{Overfitting Check: Val vs.\ Test Macro-F1}
  \label{tab:overfitting}
  \begin{tabular}{lrrl}
    \toprule
    \textbf{Model} & \textbf{Val F1} & \textbf{Test F1} & \textbf{Status} \\
    \midrule
    ViT-Small       & 99.81 & 99.65 & \textcolor{green!50!black}{OK} \\
    EfficientNet-B0 & 95.58 & 95.40 & \textcolor{green!50!black}{OK} \\
    MobileNetV3-L   & 91.44 & 91.38 & \textcolor{green!50!black}{OK} \\
    ResNet-50       & 80.11 & 80.39 & \textcolor{green!50!black}{OK} \\
    \bottomrule
  \end{tabular}
\end{table}

\subsection{Grad-CAM Interpretability}

Grad-CAM maps confirm that all three CNN models attend to agronomically
relevant image regions. For \textit{rust}, EfficientNet-B0 and ResNet-50
activate strongly on the abaxial surface regions where uredinia form,
consistent with the biological localization of the pathogen. For
\textit{miner}, activations concentrate on the serpentine gallery
boundaries, where larval damage creates the highest contrast against
healthy leaf tissue. For \textit{phoma}, EfficientNet-B0 focuses on the
necrotic center while MobileNetV3 exhibits diffuse activation around the
chlorotic halo—a behavioral difference consistent with MobileNetV3's
lower F1 on this class.

MobileNetV3 Grad-CAM required targeting the final inverted residual block
(\texttt{backbone.blocks[-1]}) rather than post-expansion layers due to
HardSwish in-place activation conflicts with PyTorch backward hooks.

ViT-Small does not support standard Grad-CAM; attention rollout or
Transformer-specific attribution methods~\cite{chefer2021} would be needed
for its interpretability analysis.

% -----------------------------------------------------------------------
\section{Discussion}

\subsection{Addressing Research Questions}

\textbf{RQ1 (CNNs vs.\ ViT):} ViT-Small outperforms all CNN architectures
by 4.25--19.26 percentage points in Macro-F1. The self-attention mechanism
enables modeling of long-range spatial dependencies across the full image,
which is advantageous for diseases with variable lesion size and
non-local spatial patterns (rust pustule distribution, miner galleries).
CNNs are limited by local receptive fields even with deep stacking,
requiring more parameters to implicitly integrate global context.

\textbf{RQ2 (Multi-source generalization):} The stratified (class $\times$
source) split proves essential. An ablation with random splitting yields
test Macro-F1 estimates 3--7 points higher on average, inflated by source
bias. With the stratified protocol, all models generalize consistently:
val--test gaps below 0.30 percentage points confirm that no split
systematically favors any source.

\textbf{RQ3 (Discriminative regions):} Grad-CAM confirms disease-specific
attention patterns for all CNN models. Clinically, the finding that
EfficientNet-B0 and ResNet-50 focus on abaxial rust pustule regions—which
agronomists use as the primary diagnostic indicator—validates that these
models have learned biologically meaningful feature representations rather
than image-level shortcuts (e.g., background color or image quality
artifacts).

\textbf{RQ4 (Accuracy--efficiency trade-off):} MobileNetV3 (5.4M
parameters, 91.38\% Macro-F1) provides the best option for resource-
constrained deployment, such as smartphone-based field applications.
EfficientNet-B0 (5.3M parameters, 95.40\%) offers superior accuracy at
identical cost when convergence time is not a constraint. ViT-Small (22.1M
parameters, 99.65\%) is optimal for cloud or server-side deployment where
accuracy is paramount.

\subsection{Practical Deployment Implications}

The three-tier architecture emerging from this study maps naturally to
deployment contexts:

\begin{enumerate}
  \item \textit{Cloud/server:} ViT-Small with API endpoint for high-stakes
    diagnostic confirmation.
  \item \textit{Edge device:} EfficientNet-B0 for tablet-based agronomist
    tools with sufficient battery and compute budget.
  \item \textit{Embedded/mobile:} MobileNetV3 for real-time smartphone
    capture applications in low-connectivity field conditions.
\end{enumerate}

\subsection{Web Application}

An interactive Streamlit application provides: (1) real-time single-image
classification with confidence scores and disease descriptions in Spanish,
(2) comparison mode running all four models simultaneously on the same
image, (3) on-demand Grad-CAM visualization for CNN models, and (4) disease-
specific treatment recommendations. The application is deployed locally with
GPU acceleration when available, with CPU fallback for inference-only
environments.

% -----------------------------------------------------------------------
\section{Threats to Validity}

\textbf{Internal validity:} All three source datasets were captured in field
or semi-controlled conditions, but image resolution, lighting, and camera
models vary substantially. Normalization with ImageNet statistics partially
mitigates this, but domain shift between datasets cannot be entirely
eliminated. Class-weighted loss reduces but does not eliminate the effect of
class imbalance (\textit{phoma} has 4{,}383 images versus 20{,}125 for
\textit{healthy}).

\textbf{External validity:} The dataset covers coffee varieties and
geographic regions specific to Brazil and Costa Rica. Disease phenotypes,
leaf morphology, and lighting conditions may differ in African or Asian
growing regions. Validation on additional field datasets (e.g., from
Ethiopia or Vietnam) is necessary before global deployment.

\textbf{Construct validity:} The five-class taxonomy is a simplification.
In practice, multiple diseases may co-occur on a single leaf, and disease
progression stages (early, mid, late) carry different agro-economic
implications. A multi-label or severity-graded formulation would better
capture operational requirements.

\textbf{Conclusion validity:} Results are consistent across 59{,}950 images
and evaluated on an independent test set with no overlap with training or
validation data. Fixed random seed (42) ensures full reproducibility.
However, the HPO budget (30 trials per model) may not fully explore the
hyperparameter space; larger budgets could further improve non-ViT models.

% -----------------------------------------------------------------------
\section{Conclusions}

This paper presented a rigorous comparative study of four deep learning
architectures for coffee leaf disease classification across five agronomic
classes. The key contributions are:

\begin{enumerate}
  \item \textbf{Unified benchmark:} 59{,}950 images from three public
    datasets, stratified by (class $\times$ source) to prevent domain leakage.
  \item \textbf{ViT superiority:} ViT-Small achieves 99.65\% Macro-F1,
    outperforming the best CNN (EfficientNet-B0, 95.40\%) by 4.25 percentage
    points, demonstrating that global self-attention is more effective than
    local convolution for this multi-disease task.
  \item \textbf{Training protocol:} Differential learning rates, label
    smoothing, class weighting, and cosine scheduling with warmup are
    essential for convergence. EfficientNet-B0 requires 50 epochs to
    saturate, emphasizing the role of training duration in compound-scaled
    architectures.
  \item \textbf{Interpretability:} Grad-CAM confirms biologically meaningful
    attention patterns in all CNN models, providing agronomic validation of
    the learned representations.
  \item \textbf{Deployment:} A Streamlit application and FastAPI endpoint
    enable real-time inference with model comparison, Grad-CAM visualization,
    and multilingual disease guidance.
\end{enumerate}

Future work includes: (1) multi-label classification for concurrent
infections, (2) disease severity grading, (3) transfer to African and
Asian coffee varieties, (4) knowledge distillation from ViT into
MobileNetV3 for accuracy--efficiency improvement, and (5) data collection
under standardized protocols to reduce source-domain variance.

% -----------------------------------------------------------------------
\section*{Technology Stack}

Python 3.11, PyTorch 2.x (CUDA 12.8, NVIDIA RTX 5060 Ti), timm 0.9,
Streamlit, FastAPI, Optuna, scikit-learn, OpenCV, Matplotlib.
All experiments use seed 42. Source code available at
\url{https://github.com/smartinezabarproyectos/CafeScan_CR}.

% -----------------------------------------------------------------------
\balance
\bibliographystyle{IEEEtran}
\bibliography{references}

\end{document}
